\section*{Chapitre \ref{ch:g7:identifying-substances} --- Identifier les espèces chimiques}

\begin{solution}{exo:g7:identifying-substances:1}
Le test à l'eau de chaux~: on fait barboter le gaz dans de l'eau de chaux
limpide, qui se trouble si le gaz est du dioxyde de carbone.
\end{solution}

\begin{solution}{exo:g7:identifying-substances:2}
La poudre est du sulfate de cuivre anhydre, et le liquide contient de
l'eau.
\end{solution}

\begin{solution}{exo:g7:identifying-substances:3}
Du dioxygène.
\end{solution}

\begin{solution}{exo:g7:identifying-substances:4}
Le crayon à papier ne \omterm{def:g2:mixing-and-dissolving:dissolve}{se dissout} pas dans le \omterm{def:g6:solutions-and-solubility:solute}{solvant}. Un trait d'encre
serait lui-même entraîné et séparé, et gâcherait le \omterm{def:g7:identifying-substances:chromatography}{chromatogramme}.
\end{solution}

\begin{solution}{exo:g7:identifying-substances:5}
Si la tache était sous le \omterm{def:g6:solutions-and-solubility:solute}{solvant}, les colorants se dissoudraient dans
le \omterm{def:g6:solutions-and-solubility:solute}{solvant} du bécher au lieu de monter le long du papier.
\end{solution}

\begin{solution}{exo:g7:identifying-substances:6}
Non, c'est un \omterm{def:g6:pure-substances-and-mixtures:mixture}{mélange} de deux colorants~: probablement le colorant jaune
et le colorant bleu des références.
\end{solution}

\begin{solution}{exo:g7:identifying-substances:7}
Recueillir un peu du gaz dans un tube à essai et approcher une bûchette
enflammée de son ouverture~: un «~pop~» aigu révèle le dihydrogène.
\end{solution}

\begin{solution}{exo:g7:identifying-substances:8}
L'\omterm{def:g4:air-a-mixture-of-gases:air}{air} expiré contient beaucoup plus de dioxyde de carbone que l'\omterm{def:g4:air-a-mixture-of-gases:air}{air}
frais.
\end{solution}

\begin{solution}{exo:g7:identifying-substances:9}
Le test montre que de l'eau est présente, pas qu'elle est seule~: de
l'eau salée, du jus ou n'importe quel \omterm{def:g6:pure-substances-and-mixtures:mixture}{mélange} aqueux bleuirait aussi la
poudre.
\end{solution}

\begin{solution}{exo:g7:identifying-substances:10}
Trois colorants (trois taches). C'est le colorant jaune qui est monté le
plus haut.
\end{solution}

\begin{solution}{exo:g7:identifying-substances:11}
Par exemple~: une bûchette incandescente dans chaque flacon --- celui
qui la rallume contient le dioxygène. Pour les deux autres, une bûchette
enflammée à l'ouverture --- un «~pop~» aigu révèle le dihydrogène. Le
dernier flacon contient le dioxyde de carbone (on peut le confirmer à
l'eau de chaux). Deux ou trois tests suffisent.
\end{solution}

\begin{solution}{exo:g7:identifying-substances:12}
Le gaz ne contient peut-être pas de dioxyde de carbone, ou trop peu pour
que cela se voie dans ce temps. Un témoin réalisé avec de l'\omterm{def:g4:air-a-mixture-of-gases:air}{air} expiré,
dont on sait qu'il contient du dioxyde de carbone, vérifie que l'eau de
chaux fonctionne.
\end{solution}

\begin{solution}{pb:g7:identifying-substances:1}
\textbf{1.} Pour que les conditions (papier, \omterm{def:g6:solutions-and-solubility:solute}{solvant}, durée) soient les
mêmes pour les quatre~: c'est seulement ainsi que l'on peut comparer les
hauteurs des taches.

\textbf{2.} Un \omterm{def:g6:pure-substances-and-mixtures:mixture}{mélange}~: son encre donne trois taches séparées, trois
colorants.

\textbf{3.} Deux colorants. Non~: l'encre du mot a trois colorants, et
le feutre 1 n'a ni le colorant rouge ni le colorant jaune.

\textbf{4.} Ses trois taches ne sont pas aux mêmes hauteurs que celles
du mot~: ce sont trois colorants différents.

\textbf{5.} Le feutre 2~: ses trois taches sont exactement aux hauteurs
des trois taches du mot.

\textbf{6.} Du dioxygène.

\textbf{7.} Du dihydrogène.

\textbf{8.} Du dioxyde de carbone.

\textbf{9.} Un témoin. Il montre que l'eau de chaux ne se trouble pas
toute seule, ni avec de l'\omterm{def:g4:air-a-mixture-of-gases:air}{air} ordinaire en ce temps~: le trouble obtenu
avec le gaz Z vient donc bien du dioxyde de carbone.

\textbf{10.} «~Eau --- sulfate de cuivre anhydre --- le blanc devient
bleu~»~: le liquide contient de l'eau.

\textbf{11.} Non. Un test positif montre que de l'eau est présente, pas
qu'il n'y a rien d'autre~: le liquide pourrait être une solution ou un
autre \omterm{def:g6:pure-substances-and-mixtures:mixture}{mélange} aqueux.

\textbf{12.} L'encre du mot contient 3 colorants, et c'est le feutre 2
qui a écrit le mot.
\end{solution}
